% Imporditud: tools/import_eksperimendid.py
% Allikas: Eesti füüsikaolümpiaadi eksperimendi ülesannete
%   kogu (Rael Kalda, 2023), ülesanne Ü46, lahendus L46.
\setAuthor{EFO žürii}
\setRound{piirkonnavoor}
\setYear{2004}
\setNumber{P/G E1}
\setDifficulty{2}
\setTopic{elektriahelad}
\setVahendid{reostaat, mõõtejoonlaud, täisnurkne kolmnurk}
\setPunktid{10}
\setAllikas{Eksperimendi kogu Ü46, lahendus lk 54}
\setLahendusseis{toores}

\prob{Reostaadi traat}
Määrata: a) reostaadile keritud traadi mass, b) reostaadi traadi eritakistus. Traadi tihedus on 8500 kg/m$^3$.

\hint

\solu
Traadi tihedus: $\rho$ = m/V , kus traadi mass m = $\rho$V (1 p). Traadi ruumala V = Sl (1 p). Traadi ristlõike pindala S = ($\pi$$d_2$)/4 (1 p). Mõõtmiste kirjeldus, kus on kirjeldatud ka kolmnurga kasutamist(1 p). Traadi pikkuse määramine (ühe keeru pikkus korda keerdude arv, kus keeru pikkus c = $\pi$d)(1 p). Keeru läbimõõdu mõõtmine: kui on arvestatud ainult keraamilise silindri või traadi kihi läbimõõtu (1 p), kui on arvestatud traadikeeru välis- ja siseläbimõõtu ning on leitud keeru läbimõõdu aritmeetiline keskmine (1 p). Traadi läbimõõdu määramine:

traadi läbimõõt = mähise pikkus ((1 p)). keerdude arv

Traadi massi arvutamine (1 p). Mõõtmistulemustest traadi eritakistuse arvutamine, reostaadi takistuse saab reostaadilt (1 p).
\probend
